\chapter{Conclusion}
\label{ch:conclusion}

\section{Ce qui a été construit}

Typhoon transforme une adresse en un enregistrement canonique de risque par
bâtiment, en joignant les aléas du référentiel Géorisques au point exact et la
vulnérabilité réelle du bâti issue de la BDNB. Sept douleurs métier ont reçu
sept réponses vérifiables, et l'application les rend utilisables dans un parcours
en quatre étapes, du diagnostic au rapport exportable.

Le résultat tient moins à une prouesse technique qu'à une série de refus :

\begin{itemize}
  \item \textbf{refus du score} — le contrat ne peut pas porter de jugement, et
    le schéma l'empêche techniquement ;
  \item \textbf{refus du zéro par défaut} — une source muette produit une absence
    déclarée, jamais une valeur ;
  \item \textbf{refus du second moteur} — le lot reste \(N\) appels au pipeline
    unitaire ;
  \item \textbf{refus de la prose non vérifiée} — le modèle propose, le
    validateur dispose.
\end{itemize}

\section{Ce que le projet démontre}

Trois enseignements nous semblent transférables au-delà de ce service.

\begin{enumerate}
  \item \textbf{La provenance est une fonctionnalité, pas une métadonnée.} Dès
    qu'un fait est vérifiable par son destinataire, la discussion change de
    nature : elle porte sur la donnée, non sur la confiance dans l'outil.
  \item \textbf{Décider de ne pas savoir est une capacité.} Distinguer cinq
    états d'absence a demandé plus de travail que d'afficher zéro — et c'est ce
    qui rend le service utilisable dans un processus d'assurance.
  \item \textbf{Un modèle de langage devient utilisable dès qu'il est encadré.}
    Non par un prompt plus ferme, mais par un validateur qui compare ce qu'il
    écrit aux faits du document et refuse le reste.
\end{enumerate}

\section{Ce qui reste ouvert}

Le rapport a nommé huit douleurs restantes plutôt que de les passer sous
silence. Les trois plus lourdes sont l'ancrage des nombres selon leur
formatage (\S\ref{ch:douleurs-restantes}, P1), la compatibilité du tier gratuit
avec des adresses réelles (P3), et l'aléa argile qui demeure communal (P5) —
cette dernière exigeant un amendement de la constitution du projet avant toute
implémentation.

\begin{invariant}[title={La règle qui tient tout le reste}]
\textbf{Aucune donnée sans source ; aucune absence confondue avec un zéro ;
aucun chiffre non ancré publié.} Le reste de l'architecture n'est qu'une
manière de rendre ces trois phrases exécutables.
\end{invariant}
